\section{Representative Case Study: Double-Word Multiplication}
\label{sec:case-study}

Double-word by double-word multiplication exercises all parts of the
architecture.  For $x=x_h+x_\ell$ and $y=y_h+y_\ell$, the algorithm computes a
leading product, accumulates cross terms, and uses a final Fast2Sum to construct
its output pair.  Its condition combines valid double-word inputs, finiteness,
and headroom for the relevant intermediate computations.

The proof separates three obligations.  First, the leading product and cross
terms must be related to the exact product $xy$.  Second, any contribution from
bounded-exponent behavior must be included in the final error statement.
Third, the final reconstruction must establish the output representation
property independently of the real-number bound.

The check proof does not repeat this analysis.  It proves instead that the C
tests imply the named multiplication condition.  At the wrapper call site,
that fact is exactly the precondition required to invoke the numerical
corollary.  This small entailment is the program-logic connection between the
runtime branch and the longer floating-point proof.

Addition provides a useful contrast.  Its proof does not impose the same lower
magnitude restrictions merely because multiplication contains an
underflow-sensitive residual.  The operation-specific architecture therefore
avoids exporting the strongest restriction of one kernel as a library-wide
policy.

For client composition, the important final step is to use only the properties
actually exported by the preceding call.  When a first call exposes the
double-word relation for its output, a second operation can consume that fact
as part of its representation reasoning; its own check remains responsible for
the other operation-specific conditions.
